\begin{definition}[Fixed-length block-injective canonical-form condition]
    \label{def:blockwise_product_word_span}
    \lean{MPSTensor.WordTupleSpanTop}
    \leanok
    For a fixed blocked length $m$, the fixed-length form of the
    block-injective canonical-form condition is
    \begin{align}
        \spn\{(A_1^\omega,\ldots,A_g^\omega):|\omega|=m\}
         & =
        \prod_{j=1}^g \MN{D_j}.
        \notag
    \end{align}
    Equivalently,
    \begin{align}
        \forall X=(X_j)_j\in\bigoplus_{j=1}^g \MN{D_j},
        \quad
        \exists (c_\omega(X))_{|\omega|=m},
        \quad
        (X_j)_{j=1}^g
         & =
        \sum_{|\omega|=m} c_\omega(X)
        (A_1^\omega,\ldots,A_g^\omega).
        \notag
    \end{align}
    This is the fixed-length form of the block-injective canonical-form
    condition of \cite[lines~2113--2117]{Cirac2021Matrix}.
\end{definition}

\begin{lemma}[BNT block separation with an injective prefix]
    \label{lem:bnt_direct_sum_injective_prefix_word_span}
    \lean{MPSTensor.hasPairBlockSeparatingWords_threeBlock_of_blocksNotGaugePhaseEquiv_c1}
    \lean{MPSTensor.propBlockInjective_of_blocksNotGaugePhaseEquiv_directSum_c1}
    \lean{MPSTensor.wordTupleSpanTop_of_blocksNotGaugePhaseEquiv_directSum_selectors_c1}
    \lean{MPSTensor.exists_wordTupleSpanTop_of_blocksNotGaugePhaseEquiv_directSum_selectors_c1}
    \leanok
    \uses{def:bnt_cpsv, def:nt_cpsv, def:blocks_not_gauge_phase_equiv,
        def:left_canonical_tensor, def:mpv_overlap,
        def:blockwise_product_word_span}
    Let $A^1,\ldots,A^r$ be normalized representatives of a basis of normal
    tensors: each representative is irreducible, the family is left-canonical,
    the self-overlaps are normalized, and no two distinct same-dimensional
    representatives are gauge-phase equivalent. Fix $L_0>0$. Assume that, for
    every block $j$, the map
    \begin{align}
        \Gamma_{L_0}^{A^j}:\MN{D_j} & \longrightarrow (\C^d)^{\otimes L_0},
        \notag                                                                     \\
        X                           & \longmapsto \sum_{|u|=L_0}\tr(XA^j_u)\ket{u}
        \notag
    \end{align}
    is injective, and that the homogeneous word spans at the two source
    lengths are full:
    \begin{align}
        \spn\{A^j_u:|u|=L_0+1\}    & = \MN{D_j},
        \notag                                   \\
        \spn\{A^j_v:|v|=3(L_0+1)\} & = \MN{D_j}.
        \notag
    \end{align}
    Then, for every ordered pair $k\ne j$, there are coefficients
    $c^{k,j}_\omega$ such that
    \begin{align}
        \sum_{|\omega|=3(L_0+1)} c^{k,j}_\omega A^k_\omega & = \Id,
        \notag                                                      \\
        \sum_{|\omega|=3(L_0+1)} c^{k,j}_\omega A^j_\omega & = 0.
        \notag
    \end{align}
    The simultaneous block-word tuples at length
    $(L_0+1)+3(r-1)(L_0+1)$ span the full product algebra
    $\prod_{j=1}^r \MN{D_j}$.
\end{lemma}

\begin{proof}
    \leanok
    \uses{def:bnt_cpsv, def:blockwise_product_word_span,
        thm:trace_nondeg, thm:cpsv_bnt_blocks_not_gaugePhaseEquiv,
        thm:normal_proportional_mpv_geometric_phase}
    Put $\ell=L_0+1$ and $S=3\ell$. For a block $A^q$, define the
    boundary-to-local-space map
    \begin{align}
        T_n^q:\MN{D_q} & \longrightarrow (\C^d)^{\otimes n},
        \notag                                               \\
        T_n^q(Z)       & =\sum_{|t|=n}\tr(ZA_t^q)\ket{t},
        \notag
    \end{align}
    and its range $\mathcal G_n^q=\operatorname{ran}T_n^q$. This is the same
    map as $\Gamma_n^{A^q}$, up to cyclically moving $Z$ inside the trace.
    The full word span at a length $n$ makes $T_n^q$ injective: if
    $T_n^q(Z)=0$, then $\tr(ZM)=0$ for every $M\in\MN{D_q}$, whence $Z=0$
    by nondegeneracy of the trace pairing.

    Fix two distinct blocks and relabel them $A$ and $B$ so that their bond
    dimensions satisfy $D_B\leq D_A$. We claim that
    \begin{align}
        \mathcal G_S^A\cap\mathcal G_S^B & =\{0\}.
        \notag
    \end{align}
    Suppose instead that $0\ne\psi$ lies in this intersection. Since the
    length-$S$ word spans are full, $T_S^A$ and $T_S^B$ are injective, so there
    are nonzero matrices $X\in\MN{D_A}$ and $Y\in\MN{D_B}$ with
    $\psi=T_S^A(X)=T_S^B(Y)$. Splitting every length-$S$ word into three
    length-$\ell$ words gives
    \begin{align}
        \tr(XA_{uvw}) & =\tr(YB_{uvw})
                      &                & (|u|=|v|=|w|=\ell).
        \notag
    \end{align}

    We now derive the local-space inclusion without suppressing the intervening
    map. Since the length-$\ell$ words of $A$ span $\MN{D_A}$ and $X\ne0$,
    the two-sided products span the same algebra:
    \begin{align}
        \spn\{A_uXA_v:|u|=|v|=\ell\} & =\MN{D_A}.
        \notag
    \end{align}
    Indeed, after expanding arbitrary left and right factors in length-$\ell$
    words, the left-hand side is the span of $MXN$ with
    $M,N\in\MN{D_A}$; a nonzero entry of $X$, together with matrix units for
    $M$ and $N$, produces every matrix unit. Hence, for any
    $Z\in\MN{D_A}$, there are scalars $a_{u,v}$ such that
    \begin{align}
        Z & =\sum_{u,v}a_{u,v}A_uXA_v.
        \notag
    \end{align}
    Define the corresponding boundary matrix for $B$ by
    \begin{align}
        R(Z) & =\sum_{u,v}a_{u,v}B_uYB_v\in\MN{D_B}.
        \notag
    \end{align}
    For each length-$\ell$ word $t$, cyclicity of the trace and the three-block
    relation above, applied to the concatenated word $vtu$, give
    \begin{align}
        \tr(ZA_t)
         & =\sum_{u,v}a_{u,v}\tr(XA_vA_tA_u)
        \notag                               \\
         & =\sum_{u,v}a_{u,v}\tr(YB_vB_tB_u)
        \notag                               \\
         & =\tr(R(Z)B_t).
        \notag
    \end{align}
    Thus $T_\ell^A(Z)=T_\ell^B(R(Z))$. Although the coefficients in a
    representation of $Z$ need not be unique, their existence for each $Z$ is
    enough to prove
    \begin{align}
        \mathcal G_\ell^A & \subseteq\mathcal G_\ell^B.
        \notag
    \end{align}

    Both $T_\ell^A$ and $T_\ell^B$ are injective, so their ranges have
    dimensions $D_A^2$ and $D_B^2$. The inclusion just proved therefore yields
    \begin{align}
        D_A^2
         & =\dim\mathcal G_\ell^A
        \leq\dim\mathcal G_\ell^B
        =D_B^2.
        \notag
    \end{align}
    On the other hand, the chosen orientation $D_B\leq D_A$ gives
    $D_B^2\leq D_A^2$. Hence $D_A=D_B$ and the inclusion is an equality.
    If the original bond dimensions were unequal, the strict reverse inequality
    would already contradict the local-space inclusion. In the remaining
    equal-dimensional case, identify the two bond index sets through this
    equality, as in the BNT same-dimension comparison. For
    $N>\ell$, let $\mathcal C_N^q(\ell)$ be the intersection of the translates
    of $\mathcal G_\ell^q$ around the periodic chain. Equality
    $\mathcal G_\ell^A=\mathcal G_\ell^B$ implies
    \begin{align}
        \mathcal C_N^A(\ell) & =\mathcal C_N^B(\ell).
        \notag
    \end{align}
    We justify the needed injective uniqueness step here. A vector satisfying
    every translated length-$\ell$ local constraint can be cut successively
    around the ring. Injectivity of $T_{L_0}^q$ makes the resulting boundary
    matrix unique at each cut, so consistency at adjacent cuts gives one
    matrix $Z$ satisfying
    \begin{align}
        ZA_i^q & =A_i^qZ \qquad (0\leq i<d).
        \notag
    \end{align}
    Since the length-$\ell$ words span $\MN{D_q}$, their commutant consists
    only of scalar matrices. Thus $Z$ is scalar, and the translated local
    constraints identify the periodic intersection, for $N>\ell$, with the
    MPS line
    \begin{align}
        \mathcal C_N^A(\ell) & =\spn\{\Psi_N(A)\},
                             &
        \mathcal C_N^B(\ell) & =\spn\{\Psi_N(B)\},
        \notag                                               \\
        \Psi_N(A)            & =\sum_{|w|=N}\tr(A_w)\ket{w},
                             &
        \Psi_N(B)            & =\sum_{|w|=N}\tr(B_w)\ket{w}.
        \notag
    \end{align}
    Consequently, equality of the local spaces forces
    $\Psi_N(A)=c_N\Psi_N(B)$ for some scalar $c_N$ at every such $N$.
    The normalized self-overlaps are the squared norms of these vectors and
    tend to one. Hence both vectors, and therefore $c_N$, are nonzero for all
    sufficiently large $N$. Theorem~\ref{thm:normal_proportional_mpv_geometric_phase}
    now turns this eventual nonzero proportionality into a gauge-phase
    equivalence. This contradicts
    Theorem~\ref{thm:cpsv_bnt_blocks_not_gaugePhaseEquiv} for the two distinct
    BNT representatives. Thus
    $\mathcal G_S^A\cap\mathcal G_S^B=\{0\}$, and hence the same assertion for
    the original ordering of the pair.

    The direct sum map
    \begin{align}
        \Phi_{k,j}:\MN{D_k}\oplus\MN{D_j}
              & \longrightarrow (\C^d)^{\otimes S},
        \notag                                      \\
        (X,Y) & \longmapsto T_S^k(X)+T_S^j(Y)
        \notag
    \end{align}
    is now injective: the trivial intersection handles the two summands, and
    each $T_S^q$ is injective. All spaces are finite-dimensional, so the dual
    map is surjective. Under the nondegenerate trace pairings, it sends word
    coefficients $(c_\omega)_\omega$ to
    \begin{align}
        \left(\sum_{|\omega|=S}c_\omega A^k_\omega,
        \sum_{|\omega|=S}c_\omega A^j_\omega\right).
        \notag
    \end{align}
    Surjectivity at $(\Id,0)$ gives the asserted pair selector.

    For fixed $k$, multiply its selectors against the other $r-1$ blocks.
    Concatenation multiplies word evaluations, so this gives coefficients
    $b_v^{(k)}$ at length $(r-1)S$ satisfying
    \begin{align}
        \sum_{|v|=(r-1)S}b_v^{(k)}A^j_v
         & =\delta_{j,k}\Id.
        \notag
    \end{align}
    Given $X_k\in\MN{D_k}$, the full length-$\ell$ word span writes
    $X_k=\sum_{|u|=\ell}a_u^{(k)}A^k_u$. Concatenating this expansion with the
    selector realizes the tuple supported at $k$ with value $X_k$, at the
    exact length
    \begin{align}
        \ell+(r-1)S
         & =(L_0+1)+3(r-1)(L_0+1).
        \notag
    \end{align}
    Summing these supported tuples over $k$ yields every element of the stated
    product algebra.
\end{proof}

\begin{theorem}[Simultaneous one-letter span after blocking a BNT]
    \label{thm:cpsv_bnt_blocked_simultaneous_span}
    \lean{MPSTensor.IsCPSVBasisOfNormalTensors.%
        exists_positive_wordTupleSpanTop}
    \lean{MPSTensor.wordTupleSpanTop_blockTensor_one}
    \lean{MPSTensor.IsCPSVBasisOfNormalTensors.%
        exists_blocked_wordTupleSpanTop_one}
    \leanok
    \uses{def:bnt_cpsv, def:blocked_tensor,
        def:blockwise_product_word_span}
    Let $(A_j)_{j=1}^g$ be a basis of normal tensors, with bond dimensions
    $D_j$. There is a positive integer $L$ such that the one-letter
    evaluations of the blocked tensors span the full product algebra:
    \begin{align}
        \spn_{\C}\left\{((A_1^{[L]})^i,\ldots,(A_g^{[L]})^i)
        : 0\leq i<d^L\right\}
         & = \prod_{j=1}^g \MN{D_j}.
        \notag
    \end{align}
    This is the existence form of simultaneous block injectivity associated
    with \cite[lines~317--345]{Cirac2016MPDO_arXiv}; no numerical bound on
    $L$ is asserted here.
\end{theorem}

\begin{proof}\leanok
    \uses{lem:cpsv_bnt_block_dimensions_positive,
        thm:cpsv_bnt_blocks_not_gaugePhaseEquiv,
        thm:cpsv_normal_tp_gauge,
        thm:normal_tp_primitive_irred,
        thm:common_exact_block_injectivity_normal_left_canonical,
        thm:wordspan_blk_inj_succ,
        thm:peripheral_primitive_irreducible_overlap_one,
        lem:gauge_invariance_block_injective,
        lem:bnt_direct_sum_injective_prefix_word_span}
    Every block has positive bond dimension, and distinct blocks are not
    gauge-phase equivalent after equal bond dimensions are identified. Put
    each block in its trace-preserving Perron gauge. The gauged blocks are
    irreducible and primitive, form a left-canonical family, and have
    normalized self-overlaps. Choose a common positive length $p$ at which
    every gauged block is block-injective. Positive-length propagation gives
    block injectivity at $p+1$ and at $3(p+1)$.

    Use the pair-selector conclusion of
    Lemma~\ref{lem:bnt_direct_sum_injective_prefix_word_span} with $L_0=p$.
    Thus, for each ordered pair $k\ne j$, there is a length-$3(p+1)$ word
    combination whose value is $\Id$ on block $k$ and $0$ on block $j$.
    For fixed $k$, multiply the $g-1$ selectors indexed by $j\ne k$.
    Concatenation multiplies their block evaluations and gives a selector of
    length
    \begin{align}
        3(g-1)(p+1)
        \notag
    \end{align}
    whose value is $\Id$ on block $k$ and $0$ on every other block. Gauge
    equivalence carries these selector identities back to the original blocks.

    The original blocks also have full length-$p$ word span, since block
    injectivity is invariant under the Perron gauge. Given
    $X_k\in\MN{D_k}$, choose a length-$p$ word combination evaluating to
    $X_k$ on block $k$, and concatenate it with the selector for $k$. This
    realizes the product-algebra tuple supported at $k$ with value $X_k$.
    Summing over $k$ gives the full simultaneous span at the constructed
    length
    \begin{align}
        L
         & =p+3(g-1)(p+1).
        \notag
    \end{align}
    Since $p>0$, this length is positive. Finally, a one-letter evaluation of
    $A_j^{[L]}$ is precisely a length-$L$ word evaluation of $A_j$. Hence
    blocking these $L$ sites turns the simultaneous word span just obtained
    into the asserted simultaneous one-letter span.
\end{proof}

Let $A$ be a tensor in canonical form, and let $A_1,\ldots,A_g$ be a basis of
normal tensors of $A$. For suitable coefficients $\mu_{j,q}$ and invertible
matrices $X_{j,q}$, write
$M_j=\diag(\mu_{j,1},\ldots,\mu_{j,r_j})$ and
$X=\bigoplus_{j=1}^{g}\bigoplus_{q=1}^{r_j}X_{j,q}$. Then
% Source: CPSV16, Section 2.3, labels eq:II_ABasicTensors and eq:II_X,
% source lines 283--301; rendered Equations (20a)--(21).
% Source-faithfulness verdict: the display labelled Eq19 in the source replaces
% direct sums by tensor products and is explicitly declared nonliteral at
% source lines 304--313; it is rendered as Equation (23).
% Index/type verdict: A^i_{alpha,beta} has physical index i and virtual indices
% alpha,beta; j and q select direct-sum summands and are not contraction legs.
\begin{align}
    A^i
     & =X\left[\bigoplus_{j=1}^{g}
             \left(M_j\otimes A_j^i\right)\right]X^{-1}
    =\bigoplus_{j=1}^{g}\bigoplus_{q=1}^{r_j}
    \mu_{j,q}X_{j,q}A_j^iX_{j,q}^{-1}.
    \label{eq:bnt_found_canonical_decomp}
\end{align}
This is \cite[Section~2.3, Equations~(20a)--(21)]{Cirac2016MPDO_arXiv}.
The decomposition (\ref{eq:bnt_found_canonical_decomp}) is used below to
separate the normal blocks.
